\chapter{Acupuncture}
\index{Acupuncture}

\section*{Definition and Overview}
Acupuncture is a therapy that originated in traditional Chinese medicine. It involves inserting very thin needles into specific points on the body, where they are left in place or gently stimulated. Traditional theory holds that the needles correct the flow of energy (``qi'') along pathways called meridians, restoring balance in the body.

From a modern scientific viewpoint, acupuncture is thought to work by stimulating nerves, releasing natural pain-relieving chemicals, and improving blood flow. It is most widely studied for pain conditions and is often used as a complementary treatment alongside conventional medicine, rather than as a replacement for it.

\section*{Epidemiology}
Acupuncture is one of the most widely used complementary therapies worldwide. in some countries, a substantial minority of adults report using it at some point, most commonly for chronic pain, arthritis, and headaches.

People who use acupuncture tend to have long-standing conditions that have not fully responded to conventional treatment, and many combine it with physiotherapy, medicines, or other therapies. Use is higher among people with chronic musculoskeletal and headache disorders, and its popularity has grown steadily as access to registered practitioners has increased.

\section*{Aetiology and Pathophysiology}
Traditional Chinese medicine explains acupuncture in terms of energy flow and balance in the body. Modern research offers a different set of explanations, several of which may act together.

\begin{itemize}
    \item \textbf{Nerve stimulation:} needles activate sensory nerves in the skin and muscle that signal the brain and spinal cord
    \item \textbf{Pain modulation:} release of endorphins and activation of the body's natural pain-relieving pathways
    \item \textbf{Local effects:} increased blood flow and changes in tissues around the needle
    \item \textbf{Placebo and context:} expectation, ritual, and attention from the practitioner can contribute to improvement
\end{itemize}

A careful interpretation of the evidence suggests that part of the benefit is due to these placebo and context effects, which does not diminish the genuine relief many people experience.

\section*{Clinical Features}
Acupuncture is not a disease, so it has no symptoms of its own. It is used to treat a range of conditions, with the strongest evidence supporting its use for certain pain problems.

\begin{itemize}
    \item \textbf{Chronic low back pain:} the best-studied use, with modest benefit for some people
    \item \textbf{Osteoarthritis:} especially of the knee, where it can reduce pain and improve function
    \item \textbf{Headaches:} including tension-type headaches and migraine
    \item \textbf{Nausea and vomiting:} after surgery or chemotherapy
    \item \textbf{Neck and shoulder pain:} some chronic cases improve
    \item \textbf{Other conditions:} less certain evidence supports its use for anxiety, insomnia, and other complaints
\end{itemize}

\section*{Differential Diagnosis}
Before relying on acupuncture, a doctor should properly assess the symptoms, because serious conditions can be mistaken for minor problems.

\begin{itemize}
    \item Back pain due to a fracture, infection, or cancer, or pressure on the spinal nerves (cauda equina syndrome) needing urgent treatment
    \item Headaches caused by high blood pressure or other treatable conditions
    \item Joint pain due to inflammatory arthritis needing disease-modifying treatment
    \item Symptoms of heart, lung, or neurological disease that require conventional care
    \item Bleeding disorders or blood-thinning medicines, which affect the safety of needling
\end{itemize}

\section*{When to Seek Medical Care}
See a doctor first to establish what is actually wrong before choosing acupuncture. Seek medical care if symptoms are severe, worsening, or accompanied by fever, weight loss, weakness, or numbness, or if they do not improve. Acupuncture should not delay the assessment of a potentially serious condition.

\textbf{Contact your local emergency services for:}
\begin{itemize}
    \item Chest pain, difficulty breathing, or sudden weakness
    \item Loss of consciousness, confusion, or a seizure
    \item A sudden severe headache, or new numbness or weakness of the face, arm, or leg
    \item Heavy bleeding or a serious injury
\end{itemize}

\section*{Diagnosis and Investigations}
Acupuncture itself requires no diagnostic tests, but a proper assessment of the condition being treated does.

\begin{itemize}
    \item \textbf{Medical assessment:} a doctor takes a history and examines you to diagnose the underlying condition
    \item \textbf{Investigations:} blood tests, X-rays, or scans such as MRI may be arranged to confirm the cause of symptoms
    \item \textbf{Practitioner assessment:} a registered acupuncturist takes their own history and assesses your general health
    \item \textbf{Safety check:} discussion of medicines, pregnancy, and any bleeding disorders before treatment begins
\end{itemize}

\section*{Management}
\begin{itemize}
    \item \textbf{Qualified practitioner:} choose a registered acupuncturist who uses sterile, single-use needles
    \item \textbf{Treatment course:} usually a course of several sessions over weeks, often weekly
    \item \textbf{Combined care:} best used alongside physiotherapy, exercise, and medicines, rather than instead of them
    \item \textbf{Exercise and self-management:} for conditions such as low back pain, staying active and strengthening remain important
    \item \textbf{Review:} treatment should be reviewed after several sessions and stopped if there is no benefit
    \item \textbf{Conventional follow-up:} continue to see your doctor and attend planned tests or reviews
\end{itemize}

\section*{Complications}
When performed by a qualified practitioner with sterile needles, acupuncture is very safe, but no treatment is completely free of risk.

\begin{itemize}
    \item Minor bruising, soreness, or small amounts of bleeding at the needle sites
    \item Dizziness or feeling faint during or after treatment
    \item Infection, mainly from unsterile equipment or poor hygiene
    \item Rarely, puncture of an internal organ, such as the lung, or injury to a nerve or blood vessel
    \item The risk of relying on acupuncture while a serious condition goes undiagnosed and untreated
\end{itemize}

\section*{Prognosis and Outcomes}
For several types of chronic pain, acupuncture produces small-to-moderate improvements for some people and is generally well tolerated. The benefit is often modest, and some of the effect appears to be due to expectation and the care provided rather than the needles alone.

Responses vary from person to person: some gain clear relief, while others notice little change. Acupuncture is not a cure for most underlying diseases, and its benefits are best seen as one part of an overall management plan built on an accurate diagnosis.

\section*{Prevention and Self-Care}
\begin{itemize}
    \item Seek a proper medical diagnosis before using acupuncture for a symptom
    \item Choose a registered practitioner who uses sterile, single-use needles
    \item Tell the practitioner about all medicines, especially blood thinners, and if you are pregnant
    \item Do not delay conventional treatment for symptoms that could be serious
    \item Review the treatment after a few sessions and stop if it is not helping
\end{itemize}

\section*{Sources and Further Reading}
The following organisations provide reliable information on acupuncture and complementary medicine.

\begin{itemize}
    \item NHS. \textit{Information on acupuncture, its uses, and safety.} https://www.nhs.uk/conditions/
    \item World Health Organization. \textit{Resources on traditional and complementary medicine.} https://www.who.int/
    \item Mayo Clinic. \textit{Guide to acupuncture and complementary therapies.} https://www.mayoclinic.org/
    \item National Center for Complementary and Integrative Health (NCCIH). \textit{Evidence summaries on acupuncture.} https://www.nccih.nih.gov/
    \item MedlinePlus. \textit{Health information on acupuncture and complementary medicine.} https://medlineplus.gov/
\end{itemize}
